\documentclass[11pt,a4paper]{article}

% Pacotes de idioma e fontes
\usepackage[utf8]{inputenc}
\usepackage[T1]{fontenc}
\usepackage[portuguese]{babel}

% Pacotes Matemáticos e Gráficos
\usepackage{amsmath}
\usepackage{amsfonts}
\usepackage{amssymb}
\usepackage{tikz} 
\usepackage{pgfplots}
\pgfplotsset{compat=1.18}

% Pacotes de Layout
\usepackage[margin=1.5cm]{geometry}
\usepackage{paracol} % Permite colunas paralelas independentes
\usepackage{enumitem}

% Definição Global de Macros Gráficas
\newcommand{\match}[4]{
    \draw[line width=1.5pt, draw=gray!80] (#1,#2) -- (#3,#4);
    \fill[black!70] (#1,#2) circle (2pt);
}

% Configuração da linha divisória e larguras das colunas
\setlength{\columnsep}{1cm} 
\setlength{\columnseprule}{0.5pt}

\begin{document}

% Cabeçalho Centralizado
\begin{center}
    \textbf{\Large PEAC-2026 | MATEMÁTICA - EAD} \\[0.2cm]
    \textbf{\large AULA 20 - Números Complexos} \\[0.2cm]
    \textit{Questões de Vestibulares de anos anteriores (UFRGS).}
\end{center}
\vspace{0.5cm}

% Início do ambiente de colunas paralelas
\begin{paracol}{2}

    % =========================================================
    % QUESTÕES
    % =========================================================

    \switchcolumn[0]*
    \textbf{1) UFRGS-2023} \\
    Considere as seguintes afirmações sobre números complexos.

    \begin{enumerate}[label=\Roman*., itemsep=0.2cm]
        \item O módulo de $z = 3 + 4i$ é $|z| = 5$.
        \item Se $u = 1 + i$ e $v = 1 - i$, então $|u \cdot v| = |u| \cdot |v|$.
        \item Para que $w = (x - 3) + (x + 4)i$ seja um número real, é necessário e suficiente que $x = 3$.
    \end{enumerate}

    Quais estão corretas?

    \begin{enumerate}[label=(\Alph*)]
        \item Apenas I.
        \item Apenas III.
        \item Apenas I e II.
        \item Apenas II e III.
        \item I, II e III.
    \end{enumerate}
    \vspace{0.4cm}

    \switchcolumn[1]
    \vspace{8cm}

    \switchcolumn[0]*
    \textbf{2) UFRGS-2018} \\
    Dados os números complexos $z_1 = (2, -1)$ e $z_2 = (3, x)$, sabe-se que $z_1 \cdot z_2 \in \mathbf{R}$. Então $x$ é igual a

    \begin{enumerate}[label=(\Alph*), itemsep=0.2cm]
        \item $-6$.
        \item $-\dfrac{3}{2}$.
        \item $0$.
        \item $\dfrac{3}{2}$.
        \item $6$.
    \end{enumerate}
    \vspace{0.4cm}

    \switchcolumn[1]
    \vspace{6.5cm}

    \switchcolumn[0]*
    \textbf{3) UFRGS-2017} \\
    Considere as seguintes afirmações sobre números complexos.

    \begin{enumerate}[label=\Roman*., itemsep=0.2cm]
        \item $(2+i)(2-i)(1+i)(1-i) = 10$.
        \item $\left(\dfrac{7}{2} + \dfrac{1}{3}i\right) + \left(\dfrac{3}{2} + \dfrac{2}{3}i\right) = \dfrac{5}{2} + \dfrac{1}{2}i$.
        \item Se o módulo do número complexo $z$ é $5$, então o módulo de $2z$ é $10$.
    \end{enumerate}

    Quais afirmações estão corretas?

    \begin{enumerate}[label=(\Alph*)]
        \item Apenas I.
        \item Apenas II.
        \item Apenas III.
        \item Apenas I e III.
        \item I, II e III.
    \end{enumerate}
    \vspace{0.4cm}

    \switchcolumn[1]
    \vspace{9.5cm}

    \switchcolumn[0]*
    \textbf{4) UFRGS-2017} \\
    As raízes do polinômio $P(x) = x^4 - 1$ são

    \begin{enumerate}[label=(\Alph*), itemsep=0.2cm]
        \item $\{i; -i; 0\}$.
        \item $\{1; -1; 0\}$.
        \item $\{1; -1; i; -i\}$.
        \item $\{i; -i; 1+i; 1-i\}$.
        \item $\{i; -i; -1+i; -1-i\}$.
    \end{enumerate}
    \vspace{0.4cm}

    \switchcolumn[1]
    \vspace{6cm}

    \switchcolumn[0]*
    \textbf{5) UFRGS-2016} \\
    Considere as igualdades abaixo.

    \begin{enumerate}[label=\Roman*., itemsep=0.2cm]
        \item $(1 - 2i)(1 + 2i) = 5$, sendo $i$ a unidade imaginária.
        \item $2^0 + 2^{-1} + 2^{-2} + 2^{-3} + \dots = 2$
        \item $1 - 2 + 3 - 4 + 5 - 6 + \dots + 99 - 100 = 50$
    \end{enumerate}

    Quais igualdades são verdadeiras?

    \begin{enumerate}[label=(\Alph*)]
        \item Apenas I.
        \item Apenas III.
        \item Apenas I e II.
        \item Apenas II e III.
        \item I, II e III.
    \end{enumerate}
    \vspace{0.4cm}

    \switchcolumn[1]
    \vspace{8.5cm}

    \switchcolumn[0]*
    \textbf{6) UFRGS-2010} \\
    Um polinômio de $5^\circ$ grau com coeficientes reais que admite os números complexos $-2+i$ e $1-2i$ como raízes, admite

    \begin{enumerate}[label=(\Alph*)]
        \item no máximo mais uma raiz complexa.
        \item $2-i$ e $-1+2i$ como raízes.
        \item uma raiz real.
        \item duas raízes reais distintas.
        \item três raízes reais distintas.
    \end{enumerate}
    \vspace{0.4cm}

    \switchcolumn[1]
    \vspace{6cm}

\end{paracol}

\vspace{1cm}

% ---------------------------------------------------------
% GABARITO
% ---------------------------------------------------------
\begin{center}
    \rule{0.6\textwidth}{0.5pt} \\[0.3cm]
    \textbf{\large GABARITO} \\[0.3cm]
    \begin{tabular}{|c|c|c|c|c|c|c|c|}
        \hline
        \textbf{1} & C & \textbf{2} & D & \textbf{3}            & D & \textbf{4} & C \\ \hline
        \textbf{5} & C & \textbf{6} & C & \multicolumn{4}{c|}{}                      \\ \hline
    \end{tabular} \\[0.3cm]
    \rule{0.6\textwidth}{0.5pt}
\end{center}

\end{document}